\newpage

% =========================================================
\chapter{Introdução}
% =========================================================

\textcolor{blue}{A introdução deve conduzir o leitor do contexto geral do problema até a
delimitação precisa da pesquisa}. Ao final do capítulo, o leitor deve
compreender o problema investigado, sua relevância, os objetivos do trabalho,
suas contribuições e a organização do documento.

% ---------------------------------------------------------
\section{Contextualização}
% ---------------------------------------------------------

O objetivo desta seção é conduzir o leitor do contexto amplo para o problema
científico ou tecnológico específico.

\begin{enumerate}
	\item Qual é o contexto científico, tecnológico, industrial ou social no qual o trabalho está inserido?
	\item Qual fenômeno, sistema, tecnologia ou problema constitui o objeto central da pesquisa?
	\item Por que esse tema é relevante atualmente?
	\item Qual problema concreto ou científico ainda precisa ser resolvido?
	\item Quais são as consequências técnicas, científicas, econômicas ou sociais desse problema?
	\item O que já foi desenvolvido para enfrentar esse problema?
	\item Quais limitações permanecem nas soluções existentes?
	\item Qual lacuna de conhecimento ou de tecnologia justifica a realização deste trabalho?
	\item Qual é exatamente a pergunta de pesquisa que o trabalho pretende responder?
	\item Qual é a hipótese, premissa ou proposição central que será investigada?
\end{enumerate}

{Resultado esperado.}
Ao final da contextualização, o leitor deve compreender
\textcolor{blue}{onde o problema está, por que ele importa e qual lacuna justifica o
	trabalho}.

% ---------------------------------------------------------
\section{Objetivos}
% ---------------------------------------------------------

\subsection{Objetivo geral}

\begin{enumerate}
	\item Qual é a finalidade científica ou tecnológica central do trabalho?
	\item Qual problema o trabalho pretende solucionar, analisar, modelar, desenvolver ou avaliar?
	\item Qual resultado concreto deverá existir ao final da pesquisa?
	\item O objetivo geral é compatível com o tempo, os recursos e o nível de um TCC?
	\item O objetivo é verificável por meio de experimentos, simulações, análise matemática ou métricas?
	\item O objetivo está diretamente relacionado à lacuna apresentada na contextualização?
\end{enumerate}

\subsubsection{Objetivos específicos}

\begin{enumerate}
	\setcounter{enumi}{6}
	\item Quais etapas são necessárias para alcançar o objetivo geral?
	\item É necessário realizar revisão bibliográfica específica?
	\item É necessário desenvolver um modelo matemático ou computacional?
	\item É necessário implementar, simular, testar, comparar ou validar uma solução?
\end{enumerate}

%Estrutura lógica recomendada.
%
%\begin{center}
%	\textbf{Problema $\rightarrow$ objetivo geral $\rightarrow$ objetivos específicos
%		$\rightarrow$ método $\rightarrow$ resultados $\rightarrow$ conclusão.}
%\end{center}
\textcolor{blue}{Note que isto será utilizado na seção de Análise e Discussões de Resultados e na seção de Conclusão
sobre seguinte princípio: O objetivo geral e os objetivos específicos aqui declarados foram cumpridos?
}

% ---------------------------------------------------------
\section{Contribuições}
% ---------------------------------------------------------

\begin{enumerate}
	\item Qual é a principal contribuição científica ou tecnológica do trabalho?
	\item O que foi efetivamente desenvolvido pelo autor?
	\item O que diferencia a solução proposta das existentes?
	\item O trabalho produz um novo modelo, método, algoritmo, arquitetura, dispositivo, sistema ou procedimento?
	\item Há contribuição teórica ou matemática?
	\item Há contribuição experimental ou computacional?
	\item Há contribuição metodológica?
	\item Há contribuição para uma aplicação prática ou industrial?
	\item Quais resultados podem ser reutilizados por outros pesquisadores?
	\item Qual conhecimento novo permanece após a conclusão do trabalho?
\end{enumerate}

\textcolor{blue}{
Pergunta-chave.
Se o trabalho fosse removido da literatura, \textbf{o que deixaria de existir?}
}

% ---------------------------------------------------------
\section{Publicações}
% ---------------------------------------------------------

\begin{enumerate}
	\item O trabalho resultou em artigo, congresso, periódico, capítulo ou outro produto científico?
	\item Qual parte do TCC originou cada publicação?
	\item Qual é a relação entre a publicação e os objetivos do TCC?
	\item Quais resultados foram apresentados na publicação?
	\item Houve diferenças entre a versão publicada e a versão apresentada no TCC?
	\item Quem são os autores da publicação?
	\item Qual é o veículo de publicação?
	\item Qual é a situação da publicação: submetida, aceita, publicada ou em preparação?
	\item Quais contribuições do TCC foram validadas ou disseminadas pela publicação?
	\item Existem direitos autorais, propriedade intelectual ou restrições associadas à publicação?
\end{enumerate}

% ---------------------------------------------------------
\section{Organização do restante do trabalho}
% ---------------------------------------------------------

\begin{enumerate}
	\item O que será apresentado no Capítulo~2?
	\item Por que o Estado da Arte é necessário para compreender o problema?
	\item O que será apresentado no Capítulo~3?
	\item Quais conceitos fundamentais precisam ser dominados pelo leitor?
	\item O que será apresentado no Capítulo~4?
	\item Como a solução proposta será formalizada?
	\item O que será apresentado no Capítulo~5?
	\item Como os resultados serão analisados e validados?
	\item Onde serão apresentadas as conclusões?
	\item Como os capítulos se conectam logicamente?
\end{enumerate}